\section{Ideal operations without HNF}
\label{sec:ideal}

In SQIsign, a left $\mathcal O_0$-ideal is represented by a rank-$4$ $\mathbb Z$-lattice in a fixed integral basis of $\mathcal O_0$. The standard implementation performs ideal arithmetic by first producing a generating set for the resulting lattice and then normalizing it via computing its HNF. While HNF is exact and canonical, computing it is not well suited to fixed-precision arithmetic: elimination steps of HNF algorithm may create intermediate coefficients that are substantially larger than those of the input generators.

In this section, we remove this source of coefficient growth. We replace the HNF algorithm by the modified LLL algorithm~\cite{Poh87,NS09}, which also gets a generating set of the lattice as an input and then outputs its basis. This does not alter the represented ideals; rather, it produces a reduced basis of the same lattice while keeping the size of intermediate integers under significantly smaller size than HNF algorithm. We then combine these bounds with the framework of~\cite{KLKL25} to derive sufficient bit sizes for the fixed-precision integer arithmetic used in SQIsign.

% In this section, we analyze ideal operations used in SQIsign.
% Then, we propose ideal operations replacing HNF algorithm with modified LLL algorithm~\cite{Poh87}, and then we analyze the bit size necessary for the fixed-precision integer arithmetic of SQIsign, based on the results of~\cite{KLKL25}.

\subsection{Modified LLL algorithm}
\label{subsec:modifiedLLL}

The modified LLL algorithm~\cite{Poh87} is particularly convenient in our setting because it takes as input an arbitrary generating family, possibly linearly dependent, and directly outputs an LLL-reduced basis of the lattice that it spans. This matches the output of ideal operations in SQIsign, which naturally produce redundant generating sets. Hence modified LLL can be used as a direct replacement for the HNF algorithm for ideal operations.
% We briefly recall the algorithm for completeness. Let $\tilde e_i$ denote the $i$th unit vector of $\mathbb R^n$, and let $v^T$ denote the transpose of a vector $v\in\mathbb R^n$. The modified LLL algorithm is given in Algorithm~\ref{alg:MLLL}.

To reduce the size bound on the algorithm, we use the modified LLL algorithm with floating-point arithmetic~\cite{NS09}, which we call \textsf{ML2} or modified LLL algorithm, for the sake of simplicity.
We briefly recall the algorithm for completeness. The \textsf{ML2} algorithm is given in Algorithm~\ref{alg:ML2}.
We also call \textsf{LLL} as the special case of \textsf{MLLL} algorithm, when $d=4$ in Algorithm~\ref{alg:ML2}.

\begin{algorithm}
\caption{\textsf{ML2}$(b_1,\dots,b_d)$~\cite[Figure~9]{NS09}}
\label{alg:ML2}
\begin{algorithmic}[1]
  \Statex \textbf{Input:} Nonzero vectors
  $(b_1,\dots,b_d)\subset \mathbb{Z}^4$ generating a rank-$4$ lattice, and
  fp-precision $\ell+1=53$.
  \Statex \textbf{Output:} An LLL-reduced basis $(b'_1,\dots,b'_4)$.

  % \State Compute exactly $G=G(b_1,\dots,b_d)$.
  \State $\bar\delta\gets \diamond(\diamond(3/4)+1)/2,\quad
  \bar r_{1,1}\gets \diamond(\langle b_1,b_1\rangle),\quad
  \kappa\gets 2,\quad \zeta\gets 0$.
  \While{$\kappa\le d$}
    \State $(b_\kappa,G,\bar r_{\kappa,*},\bar\mu_{\kappa,*},
    \bar s^{(\kappa)}_*) \gets
    \textsf{LazySizeReduce}(\kappa,\zeta,(b_1,\dots,b_d),G_{(b_1,\dots,b_d)},\bar r,\bar\mu)$. \Comment{Algorithm~\ref{alg:lazy-size-reduction} in Appendix~\ref{subsec:ML2-bound}}
    \State $\kappa'\gets \kappa$.
    \While{$\kappa\ge \zeta+2$ \textbf{and}
    $\diamond(\bar\delta\cdot \bar r_{\kappa-1,\kappa-1})
    \ge \bar s^{(\kappa')}_{\kappa-1}$}
      \State $\kappa\gets \kappa-1$.
    \EndWhile
    \For{$i=\zeta+1$ \textbf{to} $\kappa-1$}
      \State $\bar\mu_{\kappa,i}\gets \bar\mu_{\kappa',i},\quad
      \bar r_{\kappa,i}\gets \bar r_{\kappa',i}$.
    \EndFor
    \State $\bar r_{\kappa,\kappa}\gets \bar s^{(\kappa')}_\kappa$.
    \State Insert $b_{\kappa'}$ right before $b_\kappa$ and update $G$ accordingly.
    \If{$b_\kappa=0$}
      \State $\zeta\gets \zeta+1$.
    \EndIf
    \State $\kappa\gets \max(\zeta+2,\kappa+1)$.
  \EndWhile
  \State \Return $(b_{\zeta+1},\dots,b_d)$.
\end{algorithmic}
\end{algorithm}
% Modified LLL algorithm~\cite{Poh87} is an algorithm which given a generator set of a lattice, outputs an LLL-reduced basis.

% Let $\tilde e_i$ denote the $i$th unit vector of $\R^n$, and let $v^T$ denote the transpose of the vector $v\in\R^n$. Modified LLL algorithm is as follows: (Algorithm~\ref{alg:MLLL})

% \begin{algorithm}[!t]
% \caption{MLLL($a_1,\dots,a_g$)~\cite{Poh87}}
% \label{alg:MLLL}
% \begin{algorithmic}[1]
%   \Statex \textbf{Input:} Non-zero vectors $a_1, \dots, a_g$ of a lattice $\Lambda \subset \mathbb{R}^n$.
%   \Statex \textbf{Output:} A $\frac34$-LLL-reduced basis $(b_1,\cdots,b_\rho)$ of $\Lambda$.
  
%   \State $\alpha \gets 0,\ \beta \gets 0,\ \sigma \gets 0,\ \rho \gets 0,\ \tau \gets 2$
  
%   \State $\alpha \gets \alpha+1,\ \beta \gets \beta+1,\ h_\beta \gets \tilde{e}_\alpha,\ b_\beta \gets a_\alpha$, $\mu_{\beta j} \gets b_\beta^T b_j^* / B_j \ (1 \le j \le \beta-1),\ b_\beta^* \gets b_\beta - \sum_{j=1}^{\beta-1} \mu_{\beta j} b_j^*,\ B_\beta \gets (b_\beta^*)^T b_\beta^*$
%   \State \textbf{if} $B_\beta \ne 0$ \textbf{and} $\alpha < g$ \textbf{then} \textbf{go to} line 2.
%   \State \textbf{else} \textbf{set} $m \gets \tau$, $s \gets \beta$ (if $B_\beta \ne 0$), else $s \gets \beta - 1$.
  
%   \State $l \gets m-1$
  
%   \State \textbf{for} $|\mu_{ml}| > \frac{1}{2}$
%   \State \quad $t \gets \lfloor \mu_{ml} \rceil$
%   \State \quad $b_m \gets b_m - t b_l,\ h_m \gets h_m - t h_l$
%   \State \quad $\mu_{ml} \gets \mu_{ml} - t,\ \mu_{mj} \gets \mu_{mj} - t\mu_{lj} \ (1 \le j \le l-1)$
%   \State \textbf{if} $b_m = 0$ \textbf{then} \textbf{go to} line 24.
%   \State \textbf{if} $b_m \ne 0$ \textbf{and} $l < m-1$ \textbf{then} \textbf{go to} line 13.
  
%   \State \textbf{if} $B_m < \left(\frac{3}{4} - \mu_{m,m-1}^2\right) B_{m-1}$ \textbf{then} \textbf{go to} line 16.
  
%   \State $l \gets l-1$
%   \State \textbf{if} $l > 0$ \textbf{then} \textbf{go to} line 6.
%   \State \textbf{else} $m \gets m+1$. \textbf{if} $m > \beta$ \textbf{then} \textbf{terminate}, \textbf{else} \textbf{go to} line 5.
  
%   \State $\mu \gets \mu_{m,m-1},\ B \gets B_m + \mu^2 B_{m-1}$
%   \State \textbf{if} $B = 0$ \textbf{then} \textbf{go to} line 20.
%   \State \textbf{else} $\mu_{m,m-1} \gets \mu(B_{m-1}/B),\ B_m \gets B_m(B_{m-1}/B)$
%   \State \quad $\begin{pmatrix} \mu_{i,m-1} \\ \mu_{i,m} \end{pmatrix} \gets \begin{pmatrix} 1 & \mu_{m,m-1} \\ 0 & 1 \end{pmatrix} \begin{pmatrix} 0 & 1 \\ 1 & -\mu \end{pmatrix} \begin{pmatrix} \mu_{i,m-1} \\ \mu_{i,m} \end{pmatrix} \ (m+1 \le i \le \beta)$
  
%   \State $B_{m-1} \gets B$
%   \State $\begin{pmatrix} h_{m-1} \\ h_m \end{pmatrix} \gets \begin{pmatrix} h_m \\ h_{m-1} \end{pmatrix},\ \begin{pmatrix} b_{m-1} \\ b_m \end{pmatrix} \gets \begin{pmatrix} b_m \\ b_{m-1} \end{pmatrix},\ \begin{pmatrix} \mu_{m-1,j} \\ \mu_{m,j} \end{pmatrix} \gets \begin{pmatrix} \mu_{m,j} \\ \mu_{m-1,j} \end{pmatrix} \ (1 \le j \le m-2)$
%   \State \textbf{if} $m > 2$ \textbf{then} $m \gets m-1$. \textbf{go to} line 5.
  
%   \State $\sigma \gets \sigma+1,\ \beta \gets \beta-1,\ m_\sigma \gets h_m$
%   \State $b_i \gets b_{i+1},\ h_i \gets h_{i+1} \ (m \le i \le \beta)$
%   \State \textbf{if} $\alpha \ge g$ \textbf{then} \textbf{terminate}.
%   \State \textbf{else} update $b_i^*,\ \mu_{ij},\ B_i \ (m \le i \le \beta;\ 1 \le j < i)$, \textbf{set} $\tau \gets m+1$ \textbf{and} \textbf{go to} line 2.
% \end{algorithmic}
% \end{algorithm}

First, we analyze the size bound of integers in modified LLL algorithm as follows:

\begin{lemma}
\label{lem:ML2-bound}
During the execution of Algorithm~\ref{alg:ML2}, the maximum size of integers is at most $\max_{1\le i\le d} || b_i ||^2$.
\end{lemma}

\begin{proof}
See Appendix~\ref{subsec:ML2-bound}
\end{proof}
% \begin{lemma}
% \label{lem:mlll-bound}
% During the execution of Algorithm~\ref{alg:MLLL}, the size of integers is at most
% % the square of the input size.
% $\max_{1\le i\le g} ||a_i||^2$.
% \end{lemma}

% \begin{proof}
% See Appendix~\ref{subsec:mlll-bound}.
% \end{proof}

Lemma~\ref{lem:ML2-bound} shows that modified LLL algorithm does not create integers larger than the largest squared Euclidean norm of an input generator. This is precisely the feature needed for our ideal-arithmetic routines: once a given ideal operation outputs generators with controlled norm, it preserves that control. In this sense, modified LLL replaces HNF algorithm while avoiding the coefficient growth that would otherwise dominate the precision requirements.

\begin{remark}[Floating-point storage overhead]
\label{rem:ML2-fp-overhead}
Lemma~\ref{lem:ML2-bound} bounds only the size of the exact integer objects appearing in Algorithm~\ref{alg:ML2}.
However, Algorithm~\ref{alg:ML2} also stores the floating-point quantities.
Since we use fp-precision $\ell+1=53$, each floating-point quantity must be stored with a constant-size binary floating-point word supporting a $53$-bit significand.
For instance, this is exactly the significand precision of standard double precision, whose explicit fraction field has $52$ bits.
Hence the floating-point part contributes only a constant-size storage overhead per floating-point quantity, independent of the input size, but this overhead should still be counted in a concrete implementation of Algorithm~\ref{alg:ML2}.
\end{remark}

\subsection{Ideal multiplication}
\label{subsec:idealmultiplication}

We now replace the HNF step in ideal multiplication by a reduction step based on modified LLL. In the standard approach, one forms the product lattice from the sixteen pairwise products of basis elements and then computes an HNF basis of the resulting lattice. However, ideal multiplication only requires some basis of the product lattice. This makes modified LLL a natural substitute, since it accepts a redundant generating set and returns a reduced basis of the same lattice.

We propose a new algorithm for Ideal Multiplication replacing HNF algorithm with Modified LLL algorithm. The idea of the algorithm is as follows:
\begin{itemize}
\item We first clear denominators so that all generators lie in the ambient integral lattice,
\item then shorten the two input bases by LLL reduction,
\item form the sixteen products $\alpha_i\beta_j$,
\item and finally apply modified LLL to these generators. The output basis is then rescaled by the common denominator $r_1r_2$.
\end{itemize}
Our compact ideal multiplication algorithm is given in Algorithm~\ref{alg:IdealMultiplication}.

\begin{algorithm}
\caption{CompactIdealMultiplication($I_1,I_2$)}
\label{alg:IdealMultiplication}
\begin{algorithmic}[1]
  \Statex \textbf{Input:} A left $\mathcal{O}_1$-ideal $I_1 = \langle \alpha_1,\alpha_2,\alpha_3,\alpha_4 \rangle$ and a left $\mathcal{O}_2$-ideal $I_2 = \langle \beta_1,\beta_2,\beta_3,\beta_4 \rangle$ with LLL-reduced bases, such that $O_R(I_1)=O_L(I_2)$.
  \Statex \textbf{Output:} A left $\mathcal{O}_1$-ideal $I_1I_2 = \langle \gamma_1,\gamma_2,\gamma_3,\gamma_4 \rangle$.
  \State $r_1 \gets \lcm(r_{\alpha_1},r_{\alpha_2},r_{\alpha_3},r_{\alpha_4}),\
         r_2 \gets \lcm(r_{\beta_1},r_{\beta_2},r_{\beta_3},r_{\beta_4})$
  \State $\alpha_i \gets r_1\alpha_i$, $\beta_j \gets r_2\beta_j$ for $1\leq i,j \leq4$
  % \State $(\alpha_1,\alpha_2,\alpha_3,\alpha_4) \gets \textbf{Minkowski}(\alpha_1,\alpha_2,\alpha_3,\alpha_4)$
  \State $(\alpha_1,\alpha_2,\alpha_3,\alpha_4) \gets \textsf{LLL}(\alpha_1,\alpha_2,\alpha_3,\alpha_4)$
  % \State $(\beta_1,\beta_2,\beta_3,\beta_4) \gets \textbf{Minkowski}(\beta_1,\beta_2,\beta_3,\beta_4)$
  \State $(\beta_1,\beta_2,\beta_3,\beta_4) \gets \textsf{LLL}(\beta_1,\beta_2,\beta_3,\beta_4)$
  \State $M \gets (\alpha_i\beta_j)_{1\le i\le4, 1\le j\le4}$
  \State $(b_1,b_2,b_3,b_4) \gets \mathsf{MLLL}(M)$
  \State \textbf{return} $\frac{1}{r_1r_2}\langle b_1,b_2,b_3,b_4 \rangle$
\end{algorithmic}
\end{algorithm}

Correctness is immediate from the lattice interpretation of ideal multiplication. After clearing denominators, the sixteen vectors $\alpha_i\beta_j$ generate the lattice
\[
(r_1I_1)(r_2I_2)=r_1r_2(I_1I_2).
\]
Since $\mathsf{MLLL}$ returns a basis of the lattice spanned by its input vectors, the vectors $(b_1,b_2,b_3,b_4)$ form a basis of $r_1r_2(I_1I_2)$. Dividing by $r_1r_2$ therefore yields a basis of the product ideal $I_1I_2$.

Before deriving a size bound for Algorithm~\ref{alg:IdealMultiplication}, we first relate the size of the input generators to the intrinsic size of the input ideals;
in SQIsign, the quantity naturally controlled by the preceding routines is the reduced norm of an ideal, whereas an arbitrary basis of that ideal need not consist of short elements. 
For this reason, Algorithm~\ref{alg:IdealMultiplication} applies LLL reduction to the input bases before invoking $\mathsf{MLLL}$, fore the sake of bounding the (reduced) norm of each basis element depending only on $\nrd(I_1)$ and $\nrd(I_2)$.

\begin{remark}
Algorithm~\ref{alg:IdealMultiplication} needs input bases which are LLL-reduced, in order to compute the uniform bound of the maximum size of integers during the execution of the algorithm.
To compute LLL-reduced input bases, one can apply LLL algorithm to each basis of an ideal to be multiplied. The maximum size of integers during the execution of LLL algorithm is the norm of each basis element.
So, we assume that every basis element has its reduced norm less than the square of the reduced norm of the ideal and use rejection sampling for the basis of each ideal, before applying LLL algorithm.
We heuristically show that the probability that rejection sampling occurs is negligible in Appendix~\ref{subsec:exp-basis}.
\end{remark}

We now bound the largest integer that can occur during the execution of Algorithm~\ref{alg:IdealMultiplication}. The main point is that, once the input bases are LLL-reduced, the entries of the product matrix $M$ are already controlled by the norms of $I_1$ and $I_2$ by Lemma~\ref{lem:lll-norm}; Lemma~\ref{lem:ML2-bound} then propagates this control through the final $\mathsf{MLLL}$ reduction.

\begin{lemma}
\label{lem:idealmult-bound}
During the execution of Algorithm~\ref{alg:IdealMultiplication}, the size of integers is at most
\[ \frac{64p^2}{\pi^4}\, r_1^2r_2^2\,\nrd(I_1)\nrd(I_2). \]
\end{lemma}

\begin{proof}
\noindent\textbf{(1) Lines 3--4.}
% Computing Minkowski-reduced basis, the maximum size of integers is less than
% \[\frac{64p}{\pi^2}\nrd(r_tI_t)+\frac{648p}{\pi^2}\nrd(r_tI_t) = \frac{712p}{\pi^2}r_t^2\nrd(I_t), \quad(t\in\{1,2\})\]
% by Lemma~\ref{lem:minkowski-norm}, Lemma~\ref{lem:lll-norm}, and Lemma~\ref{lem:minkowski-bound}. (See Appendix~\ref{sec:minkowski}.)
Computing LLL-reduced basis, the maximum size of integers is at most
\[\max\left( \max_{1\le i\le4}\nrd(\alpha_i), \max_{1\le j\le4}\nrd(\beta_j) \right).\]

% By Minkowski's second theorem, for every $1\le i,j\le 4$,
% \[
% \nrd(\alpha_i)\le \frac{64p^2}{\pi^4}\nrd(r_1I_1),
% \qquad
% \nrd(\beta_j)\le \frac{64p^2}{\pi^4}\nrd(r_2I_2).
% \]
\smallskip
\noindent\textbf{(2) Lines 5--6.}
By Lemma~\ref{lem:lll-norm}, for every $1\le i,j\le 4$,
\[
\nrd(\alpha_i)\le \frac{8p}{\pi^2}\nrd(r_1I_1),
\qquad
\nrd(\beta_j)\le \frac{8p}{\pi^2}\nrd(r_2I_2).
\]
On the other hand, for each $1\le i,j\le4$, $\| \alpha_i\beta_j \|$ is equal to
% by~\cite[Lemma~1]{KLKL25}, every coefficient of the
% product $\alpha_i\beta_j$ in the basis $(1,i,j,k)$ is bounded by
\[
\sqrt{\nrd(\alpha_i)\nrd(\beta_j)}.
\]
Therefore every entry of the matrix $M=(\alpha_i\beta_j)_{1\le i,j\le 4}$ in
line~5 is bounded by
\[
\sqrt{\nrd(\alpha_i)\nrd(\beta_j)}
\le
% \frac{64p^2}{\pi^4}\sqrt{\nrd(r_1I_1)\nrd(r_2I_2)}.
\frac{8p}{\pi^2}\sqrt{\nrd(r_1I_1)\nrd(r_2I_2)}.
\]

By Lemma~\ref{lem:ML2-bound}, during the execution of $\mathsf{ML2}$, the maximum size is at most
\[\begin{aligned}
&\quad\max_{1\le i,j\le4}||\alpha_i\beta_j||_\infty^2
\le \max_{1\le i,j\le4}\nrd(\alpha_i\beta_j)\\
&\le \left(\frac{8p}{\pi^2}\right)^2\nrd(r_1I_1)\nrd(r_2I_2)
= \frac{64p^2}{\pi^4}\,r_1^2r_2^2\,\nrd(I_1)\nrd(I_2).
\end{aligned}\]
% By Lemma~\ref{lem:mlll-bound}, during the execution of $\mathbf{MLLL}$, the maximum size is at most
% \[\begin{aligned}
% &\quad\max_{1\le i,j\le4}||\alpha_i\beta_j||^2
% \le \max_{1\le i,j\le4}\nrd(\alpha_i\beta_j)\\
% &\le \left(\frac{64p^2}{\pi^4}\right)^2\nrd(r_1I_1)\nrd(r_2I_2)
% = \frac{2^{12}p^4}{\pi^8}r_1^2r_2^2\nrd(I_1)\nrd(I_2).
% \end{aligned}\]

\smallskip
\noindent\textbf{(3) Line 7.}
Finally, line~7 only rescales the output by the
denominator $r_1r_2$, which is smaller than the above bound. Hence, throughout
Algorithm~\ref{alg:IdealMultiplication}, the size of integers is at most
\[
\frac{64p^2}{\pi^4}\,r_1^2r_2^2\,\nrd(I_1)\nrd(I_2).
\]
\end{proof}

Lemma~\ref{lem:idealmult-bound} shows that the multiplication routine operates at the natural scale of the product ideal, up to the explicit factor $\frac{64p^2}{\pi^4}$. In particular, the reduction step based on modified LLL does not reintroduce the coefficient growth that is typical of HNF-based normalization.

\begin{remark}
In order to reduce the size bound of the norm of input bases elements, we propose computing the Minkowski-reduced basis, before computing the ideal multiplication with modified LLL algorithm.
Assuring input bases of modified LLL algorithm are LLL-reduced, we can bound their norm as $O(\nrd(I_1I_2))$, which is essentially smaller than the state-of-the-art $O(\nrd(I_1I_2)^4)$~\cite[Lemma~12]{KLKL25}.
% We provide an analysis of the algorithm for computing the Minkowski-reduced basis in Appendix~\ref{sec:minkowski}.
\end{remark}

From the result of this subsection, we also compute bound on the algorithm \textsf{IdealToIsogeny} which uses ideal multiplication as a subroutine. See Appendix~\ref{sec:idealtoisogeny}.


\subsection{Ideal intersection}
\label{subsec:idealintersection}

We also refine the integer size analysis for ideal intersection.
We also replace the HNF step by an MLLL reduction as in the case of ideal multiplication, which preserves the generated lattice up to a unimodular change of basis while keeping the intermediate coefficients smaller.

% To analyze the ideal intersection, we first analyze two subroutines used in the lattice intersection algorithm.

% \begin{algorithm}
% \caption{AdjugateWithDet($M$)}
% \label{alg:AdjugateWithDet}
% \begin{algorithmic}[1]
%   \Statex \textbf{Input:} A matrix $M\in\Z^{4\times4}$ representing a $\frac34$-LLL-reduced basis of a left $\calO_0$-ideal $I$.
%   \Statex \textbf{Output:} A pair $(A,\Delta)$ such that
%   $A=\operatorname{adj}(M)$, $\Delta=\det(M)$, and
%   $M^{-1}=A/\Delta$.

%   \State $A \gets 0_{4\times 4}$

%   \For{$i=1$ to $4$}
%     \For{$j=1$ to $4$}
%       \State $(s_1,s_2,s_3)\gets\mathsf{SORT}(\{1,2,3,4\}\setminus\{i\})$
%       \State $(c_1,c_2,c_3)\gets\mathsf{SORT}(\{1,2,3,4\}\setminus\{j\})$
%       \State Let $N^{(i,j)}=(n_{a,b})_{1\le a,b\le 3}$ be the $3\times 3$ matrix
%       defined by
%       \[
%         n_{a,b}=M_{s_a,s_b}
%         \qquad (1\le a,b\le 3).
%       \]
%       \State $\mu_{i,j} \gets
%       n_{1,1}(n_{2,2}n_{3,3}-n_{2,3}n_{3,2})
%       -n_{1,2}(n_{2,1}n_{3,3}-n_{2,3}n_{3,1})
%       +n_{1,3}(n_{2,1}n_{3,2}-n_{2,2}n_{3,1})$
%       \State $A_{j,i} \gets (-1)^{i+j}\mu_{i,j}$
%     \EndFor
%   \EndFor

%   \State $\Delta \gets
%   m_{1,1}A_{1,1}
%   +m_{1,2}A_{2,1}
%   +m_{1,3}A_{3,1}
%   +m_{1,4}A_{4,1}$

%   \If{$\Delta=0$}
%     \State \textbf{error} ``$M$ is not full rank''
%   \EndIf

%   \State \textbf{return} $(A,\Delta)$
% \end{algorithmic}
% \end{algorithm}

% \begin{lemma}
% \label{lem:adjugate-bound}
% During the execution of Algorithm~\ref{alg:AdjugateWithDet}, the maximum size of integers is less than \[\max\left(|\det M|,\,\frac{8p}{r}\,\frac{|\det M|}{\sqrt{\nrd(I)}}\right).\]
% \end{lemma}

% \begin{proof}
% See Appendix~\ref{subsubsec:adjugate-bound}
% \end{proof}

% \begin{lemma}
% \label{lem:dual-bound}
% During the execution of Algorithm~\ref{alg:LatticeDual}, the maximum size of integers is less than $\max\left(\det(M), r\cdot\max\limits_{1\le i,j \le4} M_{i,j}\right)$.
% \end{lemma}

% \begin{proof}
% See Appendix~\ref{subsubsec:dual-bound}
% \end{proof}

We now apply these routines to compute the intersection. (Algorithm~\ref{alg:LatticeIntersection})
The key difference from the HNF-based approach is Line~9: we use MLLL rather than HNF.

\begin{algorithm}
\caption{CompactLatticeIntersection($I_1,I_2$)}
\label{alg:LatticeIntersection}
\begin{algorithmic}[1]
  \Statex \textbf{Input:} Two left $\calO_0$-ideals
  $I_1=\langle \alpha_1,\alpha_2,\alpha_3,\alpha_4\rangle$ and
  $I_2=\langle \beta_1,\beta_2,\beta_3,\beta_4\rangle$.
  \Statex \textbf{Output:} The intersection lattice
  $I_1\cap I_2=\langle \gamma_1,\gamma_2,\gamma_3,\gamma_4\rangle$.
  \State $r_1 \gets \lcm(r_{\alpha_1},r_{\alpha_2},r_{\alpha_3},r_{\alpha_4})$
  \State $r_2 \gets \lcm(r_{\beta_1},r_{\beta_2},r_{\beta_3},r_{\beta_4})$
  \State $M_1 \gets r_1\cdot(\alpha_1,\alpha_2,\alpha_3,\alpha_4)$
  \State $M_2 \gets r_2\cdot(\beta_1,\beta_2,\beta_3,\beta_4)$
  \State $M_1 \gets \mathsf{LLL}(M_1;\delta=\frac34)$
  \State $M_2 \gets \mathsf{LLL}(M_2;\delta=\frac34)$
  \State $(D_1,s_1) \gets \mathsf{LatticeDual}(M_1,r_1)$ \Comment{Algorithm~\ref{alg:LatticeDual} in Appendix~\ref{subsec:latticeinter-bound}}
  \State $(D_2,s_2) \gets \mathsf{LatticeDual}(M_2,r_2)$
  \State $s \gets \lcm(s_1,s_2)$
  \State $C \gets
  \left[
    \frac{s}{s_1}D_1\ \middle|\ \frac{s}{s_2}D_2
  \right]$
  \State $H \gets \mathsf{MLLL}(C)$
    \Statex \hspace{\algorithmicindent} \Comment{$H/s$ generates $I_1^\#+I_2^\#$}
  \State $(M,r) \gets \mathsf{LatticeDual}(H,s)$
    \Statex \hspace{\algorithmicindent} \Comment{$M/r$ generates $(I_1^\#+I_2^\#)^\#$}
  % \State $M \gets \mathsf{MLLL}(M)$
  \State $(\gamma_1,\gamma_2,\gamma_3,\gamma_4)
    \gets (1\ i\ j\ k)\cdot (M/r)$
  \State $(\gamma_1,\gamma_2,\gamma_3,\gamma_4)
    \gets \mathsf{LLL}(\gamma_1,\gamma_2,\gamma_3,\gamma_4;\delta=\frac34)$
  \State \textbf{return} $\langle \gamma_1,\gamma_2,\gamma_3,\gamma_4\rangle$
\end{algorithmic}
\end{algorithm}

\begin{lemma}
\label{lem:latticeinter-bound}
During the execution of Algorithm~\ref{alg:LatticeIntersection}, the
maximum size of integers is at most
\[
  B_{\cap}
  :=
  \max\left\{
  \begin{array}{l}
    \displaystyle
    \max_{k\in\{1,2\}}
  \left(
    r_k^4\frac{p}{4}\nrd(I_k)^2,\,
    2p^2\,r_k^4\nrd(I_k)^{3/2}
  \right),\\[3mm]
    \displaystyle
    r_1^4r_2^4\frac{p^2}{16}
    \nrd(I_1)^2\nrd(I_2)^2,\\[4mm]
    \displaystyle
    16p^4\max_{k\in\{1,2\}} r_k^8\nrd(I_k)^3,\\[4mm]
    \displaystyle
    \max\left(
    \frac{p}{4}\nrd(I_1)^2\nrd(I_2)^2,
    2p^2\,\nrd(I_1)\nrd(I_2)
    \right)
  \end{array}
  \right\}.
\]
\end{lemma}

\begin{proof}
See Appendix~\ref{subsec:latticeinter-bound}.
\end{proof}

\begin{remark}
The LLL reductions in Lines~5--6 are not required for the correctness of Algorithm~\ref{alg:LatticeIntersection}.
The purpose of these reductions is instead to control the size of the subsequent dual computations.
In the previous version, \(\mathsf{LatticeDual}\) was applied directly to the denominator-cleared bases, hardening the intermediate-size estimate.
Here we first replace the input bases by \(\frac34\)-LLL-reduced bases, so that the adjugate bound for LLL-reduced \(\mathcal O_0\)-ideal bases applies before computing the duals.
\end{remark}
% \begin{lemma}
% \label{lem:latticeinter-bound}
% Let \(C_1,\ldots,C_g\) be the columns of the matrix \(C\) in Line~8, and let
% \(I_\Sigma\) denote the ideal lattice generated by \(H/s\), i.e.
% \[
%   I_\Sigma=I_1^\# + I_2^\#.
% \]
% Then, during the execution of Algorithm~\ref{alg:LatticeIntersection}, the
% maximum size of integers is at most
% \[
%   B_{\cap}
%   :=
%   \max\left\{
%   \begin{array}{l}
%     |\det(M_1M_2)|,\\[2mm]
%     \displaystyle
%     \max_{k\in\{1,2\}}
%     \left(
%       |\det M_k|,
%       8p\,\frac{|\det M_k|}{\sqrt{\nrd(I_k)}}
%     \right),\\[4mm]
%     \displaystyle
%     256p^2|\det M_1|^2|\det M_2|^2 \max\left(\frac{1}{\nrd(I_1)}, \frac{1}{\nrd(I_2)}\right),\\[3mm]
%     \displaystyle
%     |\det H|,
%     \quad
%     8p\,\frac{|\det H|}{\sqrt{\nrd(I_\Sigma)}}
%   \end{array}
%   \right\}.
% \]
% \end{lemma}

% \begin{proof}
% We first consider Lines~5--6.  For each \(k\in\{1,2\}\), the input to
% \(\mathsf{LatticeDual}\) is \(M_k/r_k\), representing the ideal lattice \(I_k\).
% By the corrected adjugate bound,
% \[
%   \|\operatorname{adj}(M_k)\|_\infty
%   \le
%   8p\,\frac{|\det M_k|}{r_k\sqrt{\nrd(I_k)}}.
% \]
% In Line~2 of Algorithm~\ref{alg:LatticeDual}, the numerator of the dual basis is
% computed as
% \[
%   D_k = r_k\operatorname{adj}(M_k)^{\mathsf T},
% \]
% up to the final gcd normalization.  Hence
% \[
%   \|D_k\|_\infty
%   \le
%   8p\,\frac{|\det M_k|}{\sqrt{\nrd(I_k)}}.
% \]
% The denominator \(s_k\) produced by \(\mathsf{LatticeDual}\) divides
% \(|\det M_k|\), and therefore the whole call in Line~\(4+k\) is bounded by
% \[
%   \max\left(
%     |\det M_k|,
%     8p\,\frac{|\det M_k|}{\sqrt{\nrd(I_k)}}
%   \right).
% \]

% Next consider Lines~7--8.  Since \(s_k\mid |\det M_k|\), we have
% \[
%   s=\operatorname{lcm}(s_1,s_2)
%   \le
%   s_1s_2
%   \le
%   |\det(M_1M_2)|.
% \]
% Thus the lcm computation itself is bounded by \(|\det(M_1M_2)|\).  The entries
% created in Line~8 are precisely the entries of the columns \(C_i\), and these are
% bounded by \(\max_i\|C_i\|\), hence by \(\max_i\|C_i\|^2\) for nonzero integral
% columns.

% Line~9 replaces the previous HNF step by an MLLL reduction.  By
% Lemma~\ref{lem:mlll-bound}, the integer growth in this step is bounded by the
% squared length of the largest input column:
% \[
%   \max_{1\le i\le g}\|C_i\|^2.
% \]
% This is the point where the HNF-based bound is improved: MLLL preserves the
% generated lattice while controlling the intermediate coefficients by the column
% lengths of \(C\).

% Finally, Line~10 calls \(\mathsf{LatticeDual}\) on the reduced basis \(H/s\), which
% generates
% \[
%   \Lambda_1^\# + \Lambda_2^\#.
% \]
% Applying the same corrected dual bound to this call gives
% \[
%   \max\left(
%     |\det H|,
%     8p\,\frac{|\det H|}{\sqrt{\nrd(I_\Sigma)}}
%   \right).
% \]
% Equivalently, the adjugate computation inside this call is bounded by
% \[
%   \max\left(
%     |\det H|,
%     8p\,\frac{|\det H|}{s\sqrt{\nrd(I_\Sigma)}}
%   \right),
% \]
% and the multiplication by \(s\) in Line~2 of
% Algorithm~\ref{alg:LatticeDual} cancels the denominator \(s\).

% The remaining lines only perform sign normalization, gcd normalization,
% coordinate reconstruction, or reduction on already bounded data, and do not
% introduce larger integers in this accounting.  Taking the maximum of the
% contributions from Lines~5--6, Lines~7--8, Line~9, and Line~10 gives the claimed
% bound \(B_{\cap}\).
% \end{proof}
% \begin{lemma}
% \label{lem:latticeinter-bound}
% During the execution of Algorithm~\ref{alg:LatticeIntersection}, the maximum size of integers is less than 
% \end{lemma}

% \begin{proof}
% By Lemma~\ref{lem:dual-bound}, the maximum size of integers in lines~5--6 is less than
% \[ \max_{k\in\{1,2\}}\left(\det(M_k), r_k\cdot\max\limits_{1\le i,j \le4} (M_k)_{i,j}\right) \enspace. \]
% By Lemma~\ref{lem:mlll-bound} and the definition of the dual lattice, the maximum size of integers in line~9 is less than
% \[\max_{1\le i\le g} r_1^2r_2^2||M_i||^2.\]
% By Lemma~\ref{lem:dual-bound}, the maximum size of integers in line~10 is less than $\det(M_1M_2)$.
% It is trivial that the maximum size of integers does not in other lines.
% So we conclude the proof.
% \end{proof}